\begin{claim}{A4.1}{Proved}{All local modules needed to distinguish degrees one and two}
Let $R=\C\{x,y,z\}$.  For $t=xy$,
\[
\Omega_{R/\C\{t\}}^1=R^3/R(y,x,0)
\simeq(x,y)\oplus R,
\]
and
\[
\Omega_{R/\C\{t\}}^2\simeq
(x,y)\oplus R/(x,y),\qquad
\Omega_{R/\C\{t\}}^3\simeq R/(x,y).
\]
For $t=xyz$,
$\Omega^1=R^3/R(yz,xz,xy)$ is torsion-free of rank two and is
not free at the triple point.  If $I=(xy,xz,yz)$, there is an exact
sequence
\[
0\to T\to\Omega^2\xrightarrow{\wedge dt} I\to0,
\qquad
T=\frac{R u\oplus R v}{(xu,\ y(u-v),\ zv)},
\quad IT=0,
\]
and $\Omega^3=R/I$.  Finally, for $t=s^m$ in
$R=\C\{s,y_1,y_2\}$,
\[
\begin{aligned}
\Omega^1&=R^2\oplus R/(s^{m-1})\,ds,\\
\Omega^2&=R\,dy_1\wedge dy_2
 \oplus\bigl(R/(s^{m-1})\bigr)^2,\\
\Omega^3&=R/(s^{m-1})\,ds\wedge dy_1\wedge dy_2.
\end{aligned}
\]
\end{claim}
\begin{proof}
The degree-one presentation is the quotient by $dt$; higher degrees
are its exterior powers.  At a double point the transverse quotient
$R^2/R(y,x)$ maps isomorphically to $(x,y)$ by
$(a,b)\mapsto xa-yb$ after an immaterial change of sign of one
basis vector.  More explicitly one may use
$(a,b)\mapsto xa-yb$ with relation $(y,x)$; its kernel is precisely
that relation, since $x,y$ are relatively prime.  In degree two,
with basis $A=dx\wedge dy,B=dx\wedge dz,C=dy\wedge dz$, the
relations are $xA=yA=0$, $yB+xC=0$.  This gives the displayed
torsion summand and the ideal $(x,y)$; wedging once more gives
$R/(x,y)$.

At the triple point the map from ordinary to logarithmic relative
one-forms sends $(a,b,c)$ to $(xa-zc,yb-zc)$.
Its kernel is precisely $R(yz,xz,xy)$: the two equations imply
$c=xyh$, $a=yzh$, $b=xzh$ by relative primality of the coordinates.
Thus it injects into $R^2$, proving torsion-freeness.  Its fibre at
the maximal ideal has dimension three although its rank is two,
so it is not free.  The same minimal-generator argument works at
the double locus.

In degree two, wedging with $dt$ is the row
$(xy,-xz,yz)$ in the basis $(A,B,C)$.
Its syzygies are generated freely by
$u=(z,y,0)$ and $v=(0,y,x)$.
To check generation, an equation $xyA-xzB+yzC=0$ implies
$x\mid C$, say $C=xC_1$; then $yA=z(B-yC_1)$ implies
$A=zA_1$, $B=y(A_1+C_1)$.  Thus the syzygy is $A_1u+C_1v$.
The three relations from $dt\wedge(dx,dy,dz)$ are
$-xu,y(u-v),zv$, proving the stated presentation of the kernel.
The ideal $I$ kills it by those relations.  Conversely the quotient
$I$ is torsion-free, so this kernel is exactly the torsion of
$\Omega^2$.  The presentation immediately gives $\Omega^3=R/I$.
All matrix identities have also been checked over $\Z[x,y,z]$ by
the supplied exact program.  For $t=s^m$ the sole degree-one
relation is $ms^{m-1}ds=0$; $m$ is a unit and taking exterior
powers gives the three formulas.
\end{proof}
